\documentclass[10pt,a4paper]{amsart}
\usepackage[margin=19mm]{geometry}
\usepackage{amsmath,amssymb,mathtools}
\usepackage[scr=rsfs,cal=euler]{mathalfa}
\usepackage{xfrac}
\usepackage[unicode,hidelinks,bookmarks=false]{hyperref}
\newcommand{\ie}{i.e.\ }
\newcommand{\cf}{cf.\ }
\newcommand{\resp}{respectively\ }
\newcommand{\Subsection}{\subsection}
\providecommand{\nobreakdash}{\nobreak\hbox{-}}
\setlength{\emergencystretch}{2em}
\multlinegap0pt
\usepackage{amssymb}
\usepackage{mathtools}
\usepackage[normalem]{ulem}
\usepackage[shortlabels]{enumitem}
\usepackage{xcolor}
\newtheorem{proposition}{Proposition}
\newcommand{\R}{\mathbb{R}}
\newcommand{\T}{\mathbb{T}}
\newcommand{\Z}{\mathbb{Z}}
\renewcommand{\epsilon}{\varepsilon}
\renewcommand{\geq}{\geqslant}
\renewcommand{\setminus}{\backslash}
\title{A structure theorem for polynomial return-time sets in minimal systems}
\author{Daniel Glasscock, Andreas Koutsogiannis, Anh N. Le, Joel Moreira, Florian K. Richter,, Donald Robertson}
\date{Source: arXiv:2511.02080v2 (CC BY 4.0)}
\begin{document}
\maketitle

\noindent\emph{Source and licence.} A structure theorem for polynomial return-time sets in minimal systems. Version arXiv:2511.02080v2 (CC BY 4.0), distributed by arXiv under the Creative Commons Attribution 4.0 International licence, http://creativecommons.org/licenses/by/4.0/, as recorded in the arXiv licence metadata for this exact version and shown on the abstract page https://arxiv.org/abs/2511.02080v2. Full licence notice: \texttt{licenses/arxiv-2511.02080-CC-BY-4.0.txt}.\par\smallskip

\noindent\textit{Editorial scope.} Counted unit: the article's proposition prop\_max\_equicont\_factor\_does\_not\_suffice (source0.tex lines 1635--1642). The article's complete original proof of it is delivered with it (source0.tex lines 1643--1700, one \texttt{proof} environment of 58 lines). No mathematics is written, repaired or re-derived: the statement and the proof are the article's own lines, in the article's own notation, and no retained line is substituted or re-flowed.  No further source-local context is needed: every cross-reference of every retained line resolves inside the two delivered spans.  Nothing was omitted from any retained span, so the package carries no omission record.  No retained line contains a citation, so the wrapper carries no bibliography and no bibliography entry.  No retained line contains a figure environment, an image-inclusion command or drawing code, so no image is delivered and the package declares no asset dependency.  Editorial formatting only, in addition: an AMS \texttt{amsart} preamble that restates the article's own declarations and macros used by the retained lines, namely the environments \texttt{proposition} on separate counters, together with the standard packages those lines need.  The retained environments print in this document as Proposition 1 in order of appearance, whereas the article prints the same items with the numbering of its own full document.  The complete native source of the article as distributed in the arXiv e-print for this exact version is bundled unchanged as \texttt{sources/188379/source0.tex}.
\begin{proposition}
\label{prop_max_equicont_factor_does_not_suffice}
There exists a minimal $2$-step nilsystem $(X, T)$ and a nonempty, open set $U \subseteq X$ such that if $\pi: X \to Z$ is the factor map to the maximal equicontinuous factor, then
\[
    R((\pi U)^\circ, (\pi U)^\circ, (\pi U)^\circ) \setminus R(U, U, U)
\]
is syndetic.
\end{proposition}

\begin{proof}
For $x \in \T = \R/\Z$, let $\lVert x \rVert$ denote the distance from $x$ to $0$.
Let $X = \T \times \T$ and $T: X \to X$ be defined by $T(x, y) = (x + \alpha, y + x)$ where $\alpha$ is an irrational number. Then
\[
    T^n(x, y) = \left(x + n \alpha, y + n x + \frac{n(n-1)}{2} \alpha\right).
\]


Fix $\epsilon > 0$ and define $U \subseteq X$ to be the ball centered at $(0, 0)$ having radius $\epsilon/4$. 
Suppose $n \in R(U, U, U)$ and let $(x_0, y_0) \in U \cap T^{-n} U \cap T^{-2n} U$. Then
\begin{align*}
        \max\{\lVert x_0 \rVert, \lVert y_0 \rVert\} &< \epsilon/4, \\
    \max\left\{\lVert x_0 + n \alpha \rVert, \left\lVert y_0 + n x_0 + \frac{n(n-1)}{2} \alpha \right\rVert\right\} &< \epsilon/4, \text{ and} \\
    \max\left\{\lVert x_0 + 2n \alpha \rVert,
    \left\lVert y_0 + 2n x_0 + \frac{2n(2n-1)}{2} \alpha \right\rVert\right\} &< \epsilon/4.
\end{align*}
It follows that
\[
    \lVert n^2 \alpha \rVert = \left\lVert \left(y_0 + 2n x_0 + \frac{2n(2n-1)}{2} \alpha\right) - 2\left(y_0 + n x_0 + \frac{n(n-1)}{2} \alpha\right) + y_0\right\rVert < \epsilon.
\]
Therefore,
\[
    R(U, U, U) \subseteq Q := \{n \in \Z: \ \lVert n^2 \alpha \rVert < \epsilon\}.
\]


The maximal equicontinuous factor of $(X, T)$ is the rotation by $\alpha$ on the first coordinate and so $\pi(U)^\circ = \pi(U)= \{x \in \T: \ \lVert x \rVert < \epsilon/4\}$. Let 
\[
    B = \{n \in \Z: \ \lVert n \alpha \rVert < \epsilon /8\}.
\]
Then for $n \in B$,
\[
    \max\{\lVert 0 \rVert, \lVert 0 + n \alpha \rVert, \lVert 0 + 2 n \alpha \rVert\} < \epsilon/4.
\]
As a result, $n \in R((\pi U)^\circ, (\pi U)^\circ, (\pi U)^\circ)$ and since $n$ is arbitrary, $B \subseteq R((\pi U)^\circ, (\pi U)^\circ, (\pi U)^\circ)$.

To show $R((\pi U)^\circ, (\pi U)^\circ, (\pi U)^\circ) \setminus R(U, U, U)$ is syndetic, it remains to show that $B \setminus Q$ is syndetic. Note that 
\[
    B \setminus Q = \{n \in \Z: \ \lVert n \alpha \rVert < \epsilon/8, \lVert n^2 \alpha \rVert \geq \epsilon\}.
\]
Consider the system $(X, S)$ with 
\[
    S(x, y) = (x + \alpha, y + 2x + \alpha)
\]
and the set
\[
    V = \{(x, y) \in X: \lVert x \rVert < \epsilon/8, \lVert y \rVert > \epsilon\}.
\]
Then $V$ is a nonempty open subset of $X$. We have
\[
    S^n(0, 0) = (n \alpha, n^2 \alpha)
\]
and therefore
\[
    B \setminus Q \supseteq \{n \in \Z: S^n (0, 0) \in V\}.
\]
Since the system $(X, S)$ is minimal, $B \setminus Q$ is syndetic. 
\end{proof}

\end{document}
